% sec_7_verification_v12.tex = v11 with the wording of Lean verification changed to "checked in Lean" / "checked by the Lean kernel" (5 places)
% and the m = 3, m = 5 entries of the adversarial search in Table tab:finite2 stated as the largest observed floating-point
% values (0.46287, 0.06390) instead of the bound "<= 0.46" (which the recorded maximum 0.46286958505344594 exceeds).
% sec_7_verification_v11.tex = v10 with the middle column of the three Lean tables 0.41 -> 0.40\textwidth
% (anonymous iacrtrans build: three 1.05pt overfull boxes); no change of content.
% v10 = v9 + \clearpage at the end so that Tables 4-9 are placed after Section 7
% (in v9 they all floated behind the references). v9 = v8 with the Lean tables placed [!htbp] (v8: Table 8 floated to the end). v8 = v7 with the last Lean table split in two (v7: float too large by 266pt in amsart).
% v7: formalization subsection rewritten for the additive two-word mixer;
% internal session labels and provenance comments removed.
\section{Finite calibration, formalization and cost: the actual scope}\label{sec:verif}

The theorems of Sections~\ref{sec:diff}--\ref{sec:mixing} are proved on paper. This section records
which parts were additionally checked on finite instances, which parts are checked by the Lean kernel,
and what was measured about cost. A finite check is a calibration, not a proof of a general statement.

\subsection{Finite calibration}
\begin{table}[p]\footnotesize
\caption{Finite checks (exact integer or rational arithmetic unless stated).}\label{tab:finite}
\begin{tabular}{p{0.25\textwidth}p{0.35\textwidth}p{0.3\textwidth}}
\toprule
statement & instances & result\\
\midrule
Theorem~\ref{thm:eps} & $m=5$: natural, $h=0$, adversarial, three random $h$; $m=7$: two $p$ with natural, $h=0$, random $h$; $m=5$ with $\wt(\rho)=3,3,5$ & $\Delta(B)=q$ in every case; $\cZ$-bound $2q/2^{\wt(s)}$ never violated\\
Theorem~\ref{thm:W2} & $m=5$, author DDT, exact matrix products & $N\cdot\EDP_2$: natural $185/128$, $h=0$ $311/128$, adversarial $31/16$, random $h$ $303/256$; $x^{-1}$ with $d_0$ (ring-polynomial witness): $159/128$; $\wt(\rho)=3$ example $75/64$ ($\le6+6/q$)\\
Theorem~\ref{thm:UE} & $m=5$: 11 tables (incl.\ $\wt(\rho)=2,3$); $m=7$: 8 tables (incl.\ $\wt(\rho)=4$) & identity of Lemma~\ref{lem:Eid} holds for every $\cG$ direction; $\mathcal E_{\cZ}=211/512$ ($m=5$), $2059/8192$ ($m=7$); $\mathcal E_B\in[3.497,3.559]$ ($m=5$), $[3.345,3.357]$ ($m=7$)\\
Section~\ref{sec:repair} ($d_0$) & power sums $m=3,5,7,9$; coefficients: 40 cases each for $m=3,5$, 8 for $m=7$; correction rule: all $\lambda$, 800 pairs; $m=5$ all $1024$ local keys for $h=a^3$ & all agree; minimum number of non-zero $c$: $27$ ($\tau=0$, bound $18$), $26$ ($\tau=1$, bound $19$)\\
Lemma~\ref{lem:norm}, Theorems~\ref{thm:top}--\ref{thm:repair}, Proposition~\ref{prop:affine} (general $d$) & own arithmetic, $m=3,4,5$; 19 maps $d$ (random $\rho$, hyperplane and $K$, plus $d_0$); coefficients for 144/144/168 keys, all $L$ and $\mu$ for $m=3$ (45\,056 keys); $\chi$ two choices; 550 affine $h$ & all agree; minimal non-zero counts $6/13/28$ ($\tau\ne\eta A$) and $7/12/28$ ($\tau=\eta A$); Boolean degree of $g_d$ exactly $2$; pass rate in $\mathcal Q_2$ ($m=3$, sample) $0.873$--$0.879$\\
Proposition~\ref{prop:pf} & author DDT $\to$ integer Walsh, $m=3$ (24\,576) and $m=5$ (7\,864\,320 components), 16 tables, $\wt(\rho)=1,2,3$ & all agree; $\cG$ entries in $\{0,\pm4\}/q^2$ only\\
\bottomrule
\end{tabular}
\end{table}
\begin{table}[p]\footnotesize
\caption{Finite checks (continued).}\label{tab:finite2}
\begin{tabular}{p{0.25\textwidth}p{0.35\textwidth}p{0.3\textwidth}}
\toprule
statement & instances & result\\
\midrule
Theorem~\ref{thm:R1} & $m=3,5,7$ (20 tables; numerical norms) & every block below its bound; $\|M_\cZ\|^2=(q/2-1)/q^2$ exactly\\
Proposition~\ref{prop:Zexact} & independent code, $m=3,5$; $d_0$ and $1\oplus K_0(x/\rho)$ ($\rho=7,11$); $h\in\{0,a^3,\text{random}\}$; 15 tables & integer Walsh of the $\cZ$ rows $=x_s$; Gram $q^2[s=t]$ exactly; SVD norm and exact Rayleigh quotient of the witness $=(q/2-1)/q^2$\\
Theorem~\ref{thm:G}, Proposition~\ref{prop:Minv} & two independent implementations; $m=3$: 16, $m=5$: 21 tables; adversarial $h$ (hill climbing); $\mathsf M$ for $p=x^{-1}$, $m=3,\dots,11$ & $g=\Psi\Sigma$ entrywise; closed form of $\mathcal B^0_z$ exact; largest observed $\|M_\cG\|^2q/\mathsf M$ (floating point): $0.46287$ ($m=3$), $0.06390$ ($m=5$); $\mathsf M\le5$ for $d_0$ type\\
Theorem~\ref{thm:steps}, Corollary~\ref{cor:inv13} & every elementary inequality for $m=5,\dots,61$ (outward rational square roots); start lemma for $m=3$, all $4095$ two-word inputs, $M=[[1,1],[1,t]]$ & all hold; $\max\|\pi_2\|_\infty=0.00464<1/64$ ($q=8$ calibrates the start lemma only)\\
Example~\ref{ex:quant}, Proposition~\ref{prop:mc} & $m=3$: all $64$ carry-mask keys and all $512\cdot8$ linear $(L,\mu)$; $m=5$: $g_d$ for $\rho=1$ and $\rho=7$ & $\ell(h)=0$, no carry-mask key vanishes, $(L,\mu)=((3,0,0),0)$ vanishes ($64$ linear pairs in all); quadratic rank modulo affine $=4=m-1$, the $(m-1)$-AND circuit agrees at all points\\
Theorem~\ref{thm:V} & decomposition: toy $m=2$, $w=2$, all $16\,777\,216$ candidates, 32 cases; $m=3$ sample of $58\,035$ & accepted sets identical\\
Theorem~\ref{thm:V}(3) & $M=5$, $Q=3$, all $120$ permutations & $\RP$ law $=\RF\mid\mathsf D$\\
\bottomrule
\end{tabular}
\end{table}

\subsection{Formalization}
The Lean developments use Lean 4.31.0-rc1 and Mathlib revision
\texttt{d568c8c0}\allowbreak\texttt{9630de09}\allowbreak\texttt{7a046763}\allowbreak\texttt{c17b9ea9}\allowbreak\texttt{9f95f950}; every listed declaration depends only on the
standard axioms \texttt{propext}, \texttt{Classical.choice}, \texttt{Quot.sound}, with no \texttt{sorry}.
\begin{table}[!htbp]\footnotesize
\caption{What is checked in Lean, and what is an input.}\label{tab:lean}
\begin{tabular}{p{0.2\textwidth}p{0.40\textwidth}p{0.3\textwidth}}
\toprule
file & checked content & not covered / input\\
\midrule
\texttt{WarmStart} & arithmetic of the warm-started recursion: from initial values $\le x$ the bound $4x^2$; $\le4/q^2$ for real $q\ge32$ & the counts of Section~\ref{sec:diff} on a concrete ring\\
\texttt{SuffixReset}, \texttt{ResetCompose} & DDT row/column sums on any finite abelian group; reset identity of Lemma~\ref{lem:reset}; two applications (four layers) & the bound on the first part is an input\\
\texttt{W2Bridge} & averaged kernel of a parallel layer with product keys is the tensor product; real averaged DP of the $w=2$ superbox equals the abstract MDS kernel; $(2,2,1,1)$, $w=2$: $\E\DP\le b_\star^2$ & the one-word bound $b_\star$ (Theorem~\ref{thm:W2}) is an input; general $w$; identification with the reference code\\
\texttt{CarryPD} & two-digit addition with an \emph{arbitrary} carry function, $B=J\circ F$, real DDT: in $\cG$ directions ($d(a+u)=d(a)+k$, $k\ne0$) $\delta_B\le q$ for all $p,h,\beta$ & $\cZ$ and $\cE$ directions; Theorem~\ref{thm:W2}; connection to \texttt{W2Bridge}\\
\bottomrule
\end{tabular}
\end{table}
\begin{table}[!htbp]\footnotesize
\caption{What is checked in Lean (continued).}\label{tab:lean2}
\begin{tabular}{p{0.2\textwidth}p{0.40\textwidth}p{0.3\textwidth}}
\toprule
file & checked content & not covered / input\\
\midrule
\texttt{CarryZ}, \texttt{CarryE\_v2} & concrete carry $\star$ (product in a fixed basis $e:\Fq\simeq\FF^m$); fibres of $x\mapsto x\star s$ have $\le q/2$ points for $s\ne0$; $\cZ$ direction $\delta_B\le q$ ($d\ne0$, $d$ two-to-one); $\cE$ direction $\delta_B\le q$ ($p$ APN along $\rho$) & exact fibre size $q/2^{\wt s}$\\
\texttt{OneRound} & all $\alpha\ne0$: $\delta_B\le q$ and $\DP\le1/q$, for $d=\varepsilon\oplus K$, $\ker K=\{0,\rho\}$, $d\ne0$, $p$ APN, every $h$ (two-to-one derived; the $\cG$ lemma is that of \texttt{CarryPD}) & ---\\
\texttt{Transport\_v2} & for any group $G$ and bijection $\varphi:G\to\Fq^2$ carrying $+$ to the carry addition: $\mathrm{ddt}=\mathrm{ddtC}$ exactly; keyed families of PD layers have $\mathrm{avgDP}\le1/q$ in the definitions of \texttt{W2Bridge} & the identification $\GR(4,m)\cong(\Fq^2,\oplus_\star)$ (Lemma~\ref{lem:add}) is the hypothesis $\varphi$; the two-block bound $4/N$\\
\texttt{RepairCore\_v2} & $A^2+B\ne0\Rightarrow$ top coefficients not both $0$; root count $\ge q-(s-1)$; for any finite $\Fq$-module $V$ and $\ell(g)=1$: surjectivity, $\ell(h+\lambda g)=\ell(h)+\lambda$, exactly one failing $\lambda$, $q\cdot\#\ker\ell=|V|$ & $\ell=\ell_d$, power sums, Lucas, $\ell_d(g_d)=1$, $\Psi=\Phi^2$\\
\texttt{ColsB\_v3}, \texttt{ColsF\_v2} & the rows of Proposition~\ref{prop:table} as counts of the real DDT of $B$ and of $F$ on concrete carry coordinates, for $\wt(\rho)=1$: $\cZ\to\cZ$ $0$, $\cZ\to\cE$ column $\le q^2/2$, $\cE\to\cZ$, $\cE\to\cE$ columns $\le2q$, $\cG\to\cZ$ $\le2$, $\cG\to\cE$ $\le4$; final layer: low digit $0$ only from $\cZ$, $\cZ$ column $\le q^2$, $\cE$ column $\le2q$, $\cG$ entries $\le2$ & general $\omega$ (only $\le1/2$ for $\cZ\to\cE$), so Theorem~\ref{thm:W2}(2)\\
\texttt{TwoStep}, \texttt{Bridge} & Theorem~\ref{thm:S} one step from the warm start and the final layer: $(T_B^2T_F)(\alpha,\gamma)\le4/N$ ($q\ge12$); the real key average of $S_2$ (with whitening) equals $T_B^2T_F$; hence \texttt{hb1}--\texttt{hb4} of \texttt{W2Bridge} are theorems and the $(2,2,1,1)$, $w=2$ average is $\le(4/N)^2$ & $\wt(\rho)=1$ only; $t\ge2$ of Theorem~\ref{thm:S}; the identification $\varphi$ of $\GR(4,m)$ with $(\Fq^2,\oplus_\star)$ remains a hypothesis\\
\texttt{RepairLink\_v2} & $\ell_d$ on actual definitions ($A_d$, $B_d$ over a transversal, Frobenius) is $\Fq$-linear; $\Delta_{g_d}=r$; repair, one failure and pass count $1-1/q$ for the actual $\ell_d$ & $\ell_d(g_d)=1$ (power sums $p_2,p_3$) is the only hypothesis; that the index set is a transversal is not stated in Lean\\
\bottomrule
\end{tabular}
\end{table}
\begin{table}[!htbp]\footnotesize
\caption{What is checked in Lean (continued): the concrete two-word mixer.}\label{tab:lean3}
\begin{tabular}{p{0.2\textwidth}p{0.40\textwidth}p{0.3\textwidth}}
\toprule
file & checked content & not covered / input\\
\midrule
\texttt{DigitSplit\_v2} & the digit split $\varphi_4:(\mathbb Z/4)^m\to\Fq^2$, $x=a+2b$ coordinatewise, carries $+$ to the carry addition (\texttt{phi4\_add}) & the ring multiplication of $\GR(4,m)$ is not used; a statement there that takes the mixer entries in the \emph{pointwise} ring $(\mathbb Z/4)^m$ has no instance (\texttt{VacuousZ4})\\
\texttt{VacuousZ4} & in the pointwise ring $(\mathbb Z/4)^m$, $m\ge1$, four unit entries and a unit determinant are incompatible (\texttt{w2\_Z4\_mixer\_hyps\_false}); hence the earlier specialization of the two-word bound to that ring is vacuous and is not counted as a verification & ---\\
\texttt{MixAdd} & two-word MDS kernel lemma over any finite abelian group $A$ for an additive automorphism $M$ of $A^2$ whose four partial maps are bijective (\texttt{mds2A\_le}); real key average of the $w=2$ superbox $=$ this kernel; $(2,2,1,1)$, $w=2$: $\E\DP\le\varepsilon^2$ (\texttt{w2A\_split\_tail\_le}) & the one-word bound $\varepsilon$ is an input\\
\texttt{ThetaZ4\_v3} & mixer $M(x,y)=(x+y,\,x+\Theta y)$ on $A=(\mathbb Z/4)^m$ for an additive automorphism $\Theta$ with $\Theta-I$ invertible; with the actual local blocks $S_2$ and $\varphi_4$: $(2,2,1,1)$, $w=2$ average $\le(4/N)^2$ (\texttt{w2\_Z4\_theta\_le}); explicit witness $m=5$, $\Theta=$ multiplication by $t$ in $(\mathbb Z/4)[t]/(t^5+t^2+1)$, with $\Theta^{-1}$ and $(\Theta-I)^{-1}$ (\texttt{theta5}, \texttt{theta5\_mixer\_mds}, \texttt{w2\_Z4\_theta5\_le}) & PD data $(p,K,\varepsilon)$ with $\wt(\rho)=1$ are hypotheses; for general $m$ the choice of $\Theta$ is on paper (Remark below)\\
\bottomrule
\end{tabular}
\end{table}
The one-round bound $\varepsilon(B)\le1/q$ of Theorem~\ref{thm:eps} is formalized for all directions on concrete
carry coordinates. For $\wt(\rho)=1$ and $w=2$, the chain Proposition~\ref{prop:table} $\to$ Theorem~\ref{thm:S} (one step) $\to$
Theorem~\ref{thm:W2}(1) $\to$ Theorem~\ref{thm:P4}(1) is checked in Lean on the real key averages on the additive group
$(\mathbb Z/4)^m$ with the digit split $\varphi_4$, for the two-word mixer $(x,y)\mapsto(x+y,\,x+\Theta y)$
(\texttt{w2\_Z4\_theta\_le}), and the mixer hypotheses are shown to be satisfiable ($m=5$). The mixer is formalized as an additive
automorphism, not as a matrix of ring scalars: in the pointwise ring $(\mathbb Z/4)^m$ a matrix with unit entries and unit
determinant does not exist (\texttt{VacuousZ4}), so the earlier specialization with such scalars proves nothing about a concrete mixer
and is superseded. That the additive group of $\GR(4,m)$ in a polynomial basis is $(\mathbb Z/4)^m$, and that multiplication by $t$ is
the companion matrix $\Theta$ of the lift $\hat f$, are used on paper; the local bound for $\wt(\rho)>1$, general $\omega$, general $w$,
Theorem~\ref{thm:UE} and Section~\ref{sec:mixing} are not formalized; of Section~\ref{sec:repair} the linear part on the actual
definitions is formalized, the power sums are not. All Lean sources needed by these declarations are released as one source
closure that builds from sources only.

\begin{remark}[The mixer for general $m$]\label{rem:theta}
Let $f\in\FF[t]$ be irreducible of degree $m\ge2$ and $\hat f$ its lift with coefficients in $\{0,1\}$. Let $\Theta$ be the matrix of
multiplication by $t$ on $(\mathbb Z/4)[t]/(\hat f)\cong(\mathbb Z/4)^m$. Then $\det\Theta=\pm\hat f(0)$ and $\det(\Theta-I)=\pm\hat f(1)$
are odd, hence units of $\mathbb Z/4$, so $\Theta$ and $\Theta-I$ are invertible and $(x,y)\mapsto(x+y,\,x+\Theta y)$ is the mixer
$\bigl[\begin{smallmatrix}1&1\\1&t\end{smallmatrix}\bigr]$ of $\GR(4,m)$ in additive form, with partial maps $I,I,I,\Theta$.
This was checked exactly for all $69$ binary irreducible $f$ with $2\le\deg f\le8$ (inverse identities), and the agreement with
multiplication by $t$ in $\GR(4,m)$ on all pairs for $m=3,5$ (two polynomials each).
\end{remark}

\subsection{Cost}
Counts are in a bit-level model with AND gates (XOR and NOT counted separately), for $m=5$. They are
reported to locate the constructions, not to claim an advantage.
\begin{table}[ht]\small
\caption{Bit-model counts ($m=5$).}\label{tab:cost}
\begin{tabular}{p{0.36\textwidth}p{0.14\textwidth}p{0.4\textwidth}}
\toprule
object & AND & averaged differential bound\\
\midrule
three-round Feistel with $x^3$, independent XOR keys & 75 & $4/N$ (\cite{AO}, Thm.~2; bijective round function, $\ge3$ rounds)\\
$S_2$ with the natural table of \cite{FukuiFFA} & 145 & $4/N$ (Theorem~\ref{thm:W2})\\
$S_2$ with an inverse-based PD table & 235 & $4/N$ (Theorem~\ref{thm:W2})\\
one $F$, PD table with $h=a^3$ & 75 (98 XOR, 1 NOT) & ---\\
correction $g_d$ alone (Proposition~\ref{prop:mc}) & $m-1$, minimal (4 for $m=5$) & --- (cost of the correction only, not of $h$)\\
\bottomrule
\end{tabular}
\end{table}
The Feistel network reaches the same constant with fewer AND gates. The differences that remain are
structural (ring differences with carry versus XOR differences; $2m$-bit ring key additions; direction-wise
structure; uniformity over all $h$; the integral structure of Section~\ref{sec:repair}), not an ordering by
cost. Ring-only and MPC costs were not measured.
\clearpage
